% ===== latex-recomposed:document =====
% flm: style file flmrecomposed.f613b89e.sty
\begin{flmFloat}{figure}{NumCap}\flmIncludeGraphics{figures/a.png}\caption{A photo}\flmLabels{\label{ref1}}\end{flmFloat}

\begin{flmFloat}{figure}{NumOnly}\flmIncludeGraphics{figures/b.png}\caption{}\flmLabels{\label{ref2}}\end{flmFloat}

\begin{flmFloat}{figure}{CapOnly}\flmIncludeGraphics{figures/c.png}\caption{A photo of \emph{C}}\phantomsection\flmLabelsWithDisplayText{Photo C}{\flmLabels{\label{ref3}}}\end{flmFloat}

\begin{flmFloat}{figure}{Bare}\flmIncludeGraphics{figures/d.png}\end{flmFloat}

See \flmRef{ref1}, \flmRef{ref2} and \flmRef{ref3}. \flmRef[S]{ref1} is a photo.
%%
% ===== latex-recomposed:style flmrecomposed.f613b89e.sty =====
\NeedsTeXFormat{LaTeX2e}
\ProvidesPackage{flmrecomposed.f613b89e}
% --- packages ---
% flm.latex.pkg.adjustbox
\RequirePackage[export]{adjustbox}
% flm.latex.pkg.caption
\RequirePackage{caption}
% flm.latex.pkg.graphicx
\RequirePackage{graphicx}
% flm.latex.pkg.hyperref
\RequirePackage{hyperref}
% flm.latex.pkg.zref-clever
\RequirePackage{zref-clever}
% --- commands and environments the library defines for the document ---
% flm.latex.def.graphics
% An image, as FLM's \includegraphics writes it:
% \flmIncludeGraphics[<keys>]{<path>}, <keys> being width=<length> and
% height=<length> where the image is shown at a set size, <path> an ordinary
% argument in which the characters # % \ are written \# \% \\.
% The image never runs wider than the line (max width, a key adjustbox adds to
% \includegraphics). Redefine the hook to show images another way.
\ProvideDocumentCommand\flmIncludeGraphics{O{} m}{\includegraphics[max width=\linewidth,#1]{#2}}
% flm.latex.def.floats
% Floats: a float of FLM's figure, table or a configured float type is
% written as \begin{flmFloat}{<type>}{<case>} <content> \caption{<caption>}
% <labels> \end{flmFloat}, <case> being NumCap (numbered, with a caption),
% NumOnly (numbered, no caption: an empty \caption{} shows the number),
% CapOnly (a caption, no number) or Bare (neither).
% \flmFloatSetUseConfig<case>{<type>}{<placement>}: how a float of the case is
%   set. By default a numbered float is set in the environment <type> at
%   <placement>, and an unnumbered one where it is written, in a center
%   environment in which \caption works as in the float <type>; a CapOnly
%   float's caption is set with no number (\caption* of the caption package).
% \flmFloatPlacementArgs<case>: the placement of a float of the case, [tbph]
%   for the numbered cases and [h] for the others.
% The content is centered. Redefine these commands to set floats differently.
\providecommand\flmFloatPlacementArgsNumCap{[tbph]}
\providecommand\flmFloatPlacementArgsNumOnly{[tbph]}
\providecommand\flmFloatPlacementArgsCapOnly{[h]}
\providecommand\flmFloatPlacementArgsBare{[h]}
\def\flm@float@infloat#1#2{\edef\flm@float@begin{\noexpand\begin{#1}#2}\def\flm@float@end{\end{#1}}}
\def\flm@float@incenter#1{\def\flm@float@begin{\begin{center}\def\@captype{#1}}\def\flm@float@end{\end{center}}}
\def\flm@float@starcaption{\let\flm@float@caption\caption\def\caption{\flm@float@caption*}}
\ProvideDocumentCommand\flmFloatSetUseConfigNumCap{m m}{\flm@float@infloat{#1}{#2}}
\ProvideDocumentCommand\flmFloatSetUseConfigNumOnly{m m}{\flm@float@infloat{#1}{#2}}
\ProvideDocumentCommand\flmFloatSetUseConfigCapOnly{m m}{\flm@float@incenter{#1}\flm@float@starcaption}
\ProvideDocumentCommand\flmFloatSetUseConfigBare{m m}{\flm@float@incenter{#1}}
\ProvideDocumentEnvironment{flmFloat}{m m}{%
  \edef\flm@float@placement{\csname flmFloatPlacementArgs#2\endcsname}%
  \csname flmFloatSetUseConfig#2\endcsname{#1}{\flm@float@placement}%
  \flm@float@begin\centering}{\flm@float@end}
% flm.latex.def.refs
% \flmRef[<options>]{<label>}: a reference that gives the name of its target
% ("Section 2"), set by \zcref of the zref-clever package.
%
% The options are the ones the document wrote, in FLM's own vocabulary. This
% definition translates each of them by its key -- a value written after the
% key is ignored, as in FLM -- and drops every key it has no translation for:
%
%   noprefix    -> noname                 the number or title alone, with no
%                                         name before it ("2", not
%                                         "Section 2")
%   S           -> S                      the first name capitalized and
%                                         unabbreviated, for a reference that
%                                         opens a sentence
%   capital     -> cap=true,abbrev=false  the names capitalized ("Lemma 1")
%   lowercase   -> cap=false,abbrev=false the names in lowercase ("lemma 1")
%   abbreviated -> abbrev=true            the abbreviated names ("Lem. 1")
%                                         where the reference type declares
%                                         them, the full names otherwise
%   anything else -> dropped              the default wording; a variant of
%                                         your own naming has no counterpart
%                                         in zref-clever, which would stop on
%                                         a key it does not know
%
% Redefine \flmRef to set references differently -- another package, another
% wording, your own variants, no link. Your definition sees the options as the
% document wrote them, in FLM's vocabulary, and maps them as it likes.
%
% The translation uses expl3, LaTeX's programming layer, which is part of the
% format and needs no package; its names read only between \ExplSyntaxOn and
% \ExplSyntaxOff. \keyval_parse:NNn reads the option list as LaTeX reads any
% key-value list: it hands a key written alone to the first macro below and a
% key with its value to the second, which drops the value. \str_case:nnF looks
% the key up in the table and appends what \zcref is to be given, or nothing.
% The macros with @ in their names are internals, not a customization point:
% redefine \flmRef, not its parts.
\ExplSyntaxOn
\def\flm@ref@add#1{\tl_put_right:Nn \flm@ref@options {#1,}}
\def\flm@ref@key#1{\str_case:nnF {#1}
  {{noprefix}{\flm@ref@add{noname}}
   {S}{\flm@ref@add{S}}
   {capital}{\flm@ref@add{cap=true,abbrev=false}}
   {lowercase}{\flm@ref@add{cap=false,abbrev=false}}
   {abbreviated}{\flm@ref@add{abbrev=true}}}
  {}}
\def\flm@ref@pair#1#2{\flm@ref@key{#1}}
\def\flm@ref@zcref#1#2{\zcref[#1]{#2}}
\ProvideDocumentCommand\flmRef{O{} m}{%
  \def\flm@ref@options{}%
  \keyval_parse:NNn \flm@ref@key \flm@ref@pair {#1}%
  \exp_args:No \flm@ref@zcref {\flm@ref@options} {#2}}
\ExplSyntaxOff
% flm.latex.def.label-display-text
% \flmLabelsWithDisplayText{<text>}{<labels>}: makes <labels> refer to this
% place, with <text> as what a reference to them prints. It is the only
% command for that: a heading the class did not number uses it, and so does
% anything else that has a name but no number.
%
% It places no anchor of its own. Under hyperref a sectioning command, starred
% or not, already places one at the title; a \label after it points there.
% \refstepcounter would place a second anchor here, below the title, and a
% link would lead below the heading. So the counter is stepped with
% the plain \stepcounter and the two things this needs are set here:
% \@currentlabel, which is what a \label records, and \@currentcounter,
% which is what zref-clever reads to give the label its reference type. The
% group keeps every change local.
%
% <text> is expanded to make \@currentlabel, and \label does not survive that
% expansion; so a \label inside <text> (for example a footnote's) is dropped
% from the label's text only. At its own place in the document it still makes
% its own label.
%
% A caller with no sectioning command before it -- for example a defined term
% in running text -- writes \phantomsection at its own place, so that a link
% has an anchor of its own to lead to.
%
% Redefine \flmLabelsWithDisplayText to change how such labels are made.
\ifcsname c@flm@internal@refcnt\endcsname\else\newcounter{flm@internal@refcnt}\fi
\ProvideDocumentCommand\flmLabelsWithDisplayText{m m}{%
  \begingroup
    \stepcounter{flm@internal@refcnt}%
    {\let\label\@gobble\protected@xdef\flm@lwdt@text{#1}}%
    \let\@currentlabel\flm@lwdt@text
    \def\@currentcounter{flm@internal@refcnt}%
    #2%
  \endgroup}
% flm.latex.def.zref-clever-type
% The zref-clever reference type of the counter \flmLabelsWithDisplayText
% steps, with every name prefix empty, so that a reference to such a label
% prints its text alone ("Introduction", not "Section Introduction").
% Every other label keeps the name its own type gives it. Redefine these
% settings, or the type of the counter, to put something before the text.
\zcsetup{countertype={flm@internal@refcnt=flm--internal--ref}}
\zcRefTypeSetup{flm--internal--ref}{Name-sg = , name-sg = , Name-pl = , name-pl = ,}
% flm.latex.def.labels
% \flmLabels{<\label commands>}: the labels of one place in the document.
% In text it runs them as written. On a line of a display-math
% environment, where amsmath keeps the last label of a line and loses the
% ones before it, it runs the first and records every further one as an alias:
% a line per alias is written to the .aux file after the last page and read on
% the next run, making the alias refer to what the first label refers to
% (for \ref and for \zcref alike). Redefine \flmLabels to change what a
% group of labels does; \def\flmLabels#1{#1} restores plain LaTeX behavior,
% amsmath's "Multiple \label's" message included.
\ProvideDocumentCommand\flmLabels{m}{%
  \ifmmode\expandafter\flm@labels@math\else\expandafter\@firstofone\fi{#1}}
% Reading the group: the first \label as written, the rest as aliases of it.
\def\flm@labels@math#1{\def\flm@lbl@first{}\flm@lbl@scan#1\flm@lbl@stop}
\def\flm@lbl@scan{\@ifnextchar\label\flm@lbl@take\flm@lbl@rest}
\def\flm@lbl@take\label#1{%
  \ifx\flm@lbl@first\@empty
    \def\flm@lbl@first{#1}%
    \label{#1}%
  \else
    \expandafter\flm@recordalias\expandafter{\flm@lbl@first}{#1}%
  \fi
  \flm@lbl@scan}
\def\flm@lbl@rest#1\flm@lbl@stop{#1}
% Collecting the aliases. A name is recorded once however often the line is
% typeset (amsmath measures a line before setting it).
\def\flm@aliaslist{}
\def\flm@recordalias#1#2{%
  \ifcsname flm@alias@#2\endcsname\else
    \expandafter\gdef\csname flm@alias@#2\endcsname{}%
    \g@addto@macro\flm@aliaslist{\flm@alias@entry{#2}{#1}}%
  \fi}
% Writing them, from the hook that runs after the last page has been shipped
% out: every \newlabel of that page has reached the .aux file by then, so the
% alias lines come after them. \AtEndDocument would run too early.
\def\flm@alias@write#1#2{\immediate\write\@auxout{\string\flm@aliaslabel{#1}{#2}}}
\AddToHook{enddocument/afterlastpage}{\let\flm@alias@entry\flm@alias@write\flm@aliaslist}
% Reading them back: the alias refers to what the original refers to. The
% assignments are global because the .aux file is read inside a group.
\def\flm@aliaslabel#1#2{%
  \ifcsname r@#2\endcsname
    \global\expandafter\let\csname r@#1\expandafter\endcsname\csname r@#2\endcsname
  \fi
  \ifcsname Z@R@#2\endcsname
    \global\expandafter\let\csname Z@R@#1\expandafter\endcsname\csname Z@R@#2\endcsname
  \fi}
% --- definitions written from the configuration of FLM features ---
% flm@float@reftype@figure (the floats feature's, written from its configuration): declared unless \flm@float@reftype@figure is already defined, so a definition made earlier under that name is kept
\ifcsname flm@float@reftype@figure\endcsname\else\zcRefTypeSetup{figure}{Name-sg={Figure},name-sg={Fig.},Name-pl={Figures},name-pl={Figs.}}\def\flm@float@reftype@figure{}\fi
\endinput
